\section{Research log}

\subsection*{Round 0 --- private workbench initialization}

\begin{itemize}[leftmargin=1.8em]
\item Preserved the completed 121-page annotated audit as \Source{audit_baseline/annotated_audit.tex} and the 108-page reconstruction as \Source{audit_baseline/clean_transcription.tex}.
\item Converted the review objections into named mathematical obligations with explicit maps, bases, signs, and falsification tests.
\item Proved the exact conditional endpoint: once the geometric maps are identified with the printed matrices, the verified algebra gives the standard smooth $6$-sphere.
\item Kept the conductor mechanism separate from both the existence of $X$ and the truth or falsity of the published CDP theorem.
\end{itemize}

\subsection*{Next calculations}

\begin{enumerate}[leftmargin=1.8em]
\item Extract the actual coordinate formulas for the three collar inclusions and build a complete orientation/sign table.
\item Build a finite integral chain model for the cusp central fibre $W$ and compare its specialization map with Appendix~B in primitive bases.
\item Compute the Leray $d_2$ geometrically and compare it to $d_2(12\gamma)=\pm p\omega$ without importing the Mayer--Vietoris normalization.
\item Compute the restriction map
$r=\eta^*i^*:\operatorname{Pic}(X)\to\operatorname{Pic}(\mathrm{dP}_6)$
from transition functions, and complete the line-by-line CDP cotangent-sheaf
comparison without using the unresolved topology.
\end{enumerate}

\subsection*{Round 1 --- 26 August 2026: typed workbench release}

\begin{itemize}[leftmargin=1.8em]
\item Rebuilt the workbench as a new clean tree with the exact title
``S6 Proof Repair and Disproof Workbench --- Private''.  The public audit
project is not an input to the upload workflow and remains untouched.
\item Kept the terminal algebra \Verified{}, but kept the global topology
\Conditional{} until C1--C3, V1, MV1, and L1 identify the printed matrices with
actual geometric maps.  No theorem-level CDP verdict is inferred from a
proof-step defect alone.
\item Added a reader-facing literature audit with DOI-resolvable references
for CDP20/CDP98, analytic duality and base change, nearby cycles, topology,
and the Leiden Declaration.  Local SGA TeX was used as a comparison layer;
third-party files are not bundled.
\item Added the bounded Fable/Jacobian provenance question to the
comparative workstream.  Any source path not independently located remains
\Open{} rather than being presented as a recovered transcript.
\end{itemize}

\subsection*{Round 2 --- 26 August 2026: duality and base-change bridges}

\begin{itemize}[leftmargin=1.8em]
\item Proved the abstract compact-Cartier-fibre duality statement from
Ramis--Ruget's analytic dualizing complex and trace pairing, Cartier
adjunction, and coherent-cohomology finiteness.  Normality of the fibre is
not required.
\item Promoted the already complete normalization--conductor calculation to
\Verified{} and separated it from its \Conditional{} application to the
candidate.  Also isolated the elementary $\mathrm{dP}_6$ Picard pullback
lemma: $H^2(X;\mathbb Z)=0$ forces every line bundle from $X$ to pull back
trivially to the normalization.
\item At this round, promoted only the abstract theorem to \Verified{} and
left its candidate application \Conditional{} on smoothness, compactness,
the full scheme-theoretic fibre, and Cartierness.  Round~3 below supersedes
that interim status after independently retracing Sections~3--6.
\item Proved the required abstract top-degree base-change isomorphism from
the degree-$3$ vanishing corollary of the analytic base-change theorem.
Also proved that properness, flatness, Cartierness of the full fibre, and
$f$-flatness of a vector bundle are automatic for a nonconstant map from a
compact connected smooth complex manifold to a smooth curve.
\item At this round, kept the candidate instantiation \Open{} pending a chart
audit.  Round~3 below records the completed source reconstruction and status
promotion.
\item Reconciled the conditional topology endpoint with the six exact bridge
obligations C1--C3, V1, MV1, and L1; no four-item shorthand now substitutes
for that dependency set.
\end{itemize}

\subsection*{Round 3 --- 26 August 2026: candidate geometry and the
$\operatorname{Pic}^0$ strengthening}

\begin{itemize}[leftmargin=1.8em]
\item Independently retraced Theorem~3.4, Theorem~4.5,
Proposition~4.6, Proposition~4.7(i)--(iii), Theorem~5.4, and
Theorem~6.2 in the frozen transcription.  This closes the analytic action,
quotient, global gluing, full-fibre, $\mathrm{dP}_6$ normalization, and local
canonical-triviality gates.  It does not close the cusp-collapse or integral
topology maps.
\item Proved a topology-free Picard lemma: for every
$L\in\operatorname{Pic}^0(X)$, the line bundle
$A_L=(L^\vee\otimes\omega_X)|_W$ pulls back trivially to
$\mathrm{dP}_6$.  The proof uses $c_1(L)=0$, injectivity of the Chern-class
map on $\mathrm{dP}_6$, and the nowhere-vanishing three-form on $N_0$.
\item Combined that lemma with the verified conductor section, analytic
duality, and top-degree base change to obtain
$R^2f_*(T_X\otimes L)\ne0$ for every
$L\in\operatorname{Pic}^0(X)$.  This conclusion no longer borrows
$H^2(X;\mathbb Z)=0$.
\item Isolated the exact remaining all-line-bundle question as
$r=\eta^*i^*=0$.  The open global topology would imply it, but the
$\operatorname{Pic}^0$ conclusion needed for generic CDP twists is already
independent of that route.
\end{itemize}
